% Abstract, \input from main.tex
\chapter*{Abstract}
%% FUNCTION: States the question, the method, the main findings and the recommendation in one
%% paragraph.

Round Robin, in which people take turns picking their favourite remaining good, is a fair
division procedure that can be carried out by hand. It always guarantees envy-freeness up to one
good (EF1), but it offers no guarantee, not even an approximate one, for the stronger notion of
envy-freeness up to any good (EFX). We therefore ask how often Round Robin achieves EFX in the
average case and under which conditions it is a good heuristic for EFX. We simulate
$20{,}000$ instances for each combination of two to five agents, up to four goods per agent and
lognormal values with five levels of agreement and six levels of dispersion, and vary one
condition at a time around a reference cell. In the reference cell with four agents and twelve
goods, Round Robin achieves EFX in $29\%$ of the instances. It achieves EFX more often when agents
agree little, when values are similar and the number of goods is divisible by the number of
agents, and when every agent receives more goods; it achieves EFX less often, and misses it by
more, when agents agree strongly, values differ greatly or more agents take part. We prove that
Round Robin is always EFX when every agent receives at most two goods and all favourite goods
differ, and that it is never EFX under identical valuations when the number of goods is not
divisible by the number of agents. Compared with a random allocation of equal bundle sizes,
Round Robin is far ahead unless agents agree strongly, and reversing the order of the last round
removes its systematic failure under identical valuations. A decision table summarises when Round
Robin can be relied on.
